\documentclass[letterpaper,twocolumn,10pt]{article}
\usepackage[margin=0.75in]{geometry}
\usepackage{times}
\usepackage{booktabs}
\usepackage{xspace}
\usepackage{enumitem}
\usepackage{listings}
\usepackage{amsmath}
\usepackage{url}
\usepackage{graphicx}
\setlength{\emergencystretch}{2em}

\newcommand{\sys}{ContractBPF-Ledger\xspace}

\title{\bf ContractBPF-Ledger: Bounded Resource-Effect Accounting for BPF-Programmable Scheduling and Paging}
\author{Paper ID XXX\\Anonymous Institution}
\date{}

\begin{document}
\maketitle

\begin{abstract}
Networked services increasingly rely on programmable kernel policies to meet workload-specific latency, throughput, and resource-efficiency goals. eBPF and related mechanisms now make it possible to specialize core policy paths such as scheduling and memory management. Existing verifier, isolation, and language-based mechanisms protect important local safety properties, including memory safety, bounded control flow, helper-call validity, and pointer safety. They do not directly reason about the dynamic resource effects of multiple independently loaded policies. A BPF scheduler and a BPF paging policy may both pass verification and improve their own objectives in isolation, yet their composition can create refault storms, queue-delay amplification, and tail-latency collapse for a multi-tenant networked service.

This paper presents \sys, a resource-effect accounting system for BPF-programmable scheduling and paging policies. \sys grants policies effect tokens, records their resource effects in per-scope ledgers, and enforces budgets at effect boundaries rather than at every BPF instruction. When a policy exceeds its budget or triggers a cross-subsystem feedback loop, \sys performs bounded degradation: it throttles or revokes the smallest harmful effect, such as page demotion, while preserving unrelated safe policy functionality.

We implement \sys on Linux using sched\_ext and a conservative PageFlex-style paging decision hook. The current artifact validates the mechanism in QEMU: a controlled G1--G9 experiment matrix reproduces the scheduler-paging feedback loop and shows recovery by revoking only demote\_page, while a memcached G1--G9 companion matrix demonstrates the same artifact path on a real networked service. These QEMU results are correctness and reproducibility evidence, not production performance claims.
\end{abstract}

% NSDI '27 prescreening note: keep the Introduction within three pages after the abstract.
\section{Introduction}

Cloud and multi-tenant networked services are increasingly sensitive to kernel resource policy. A small shift in scheduling, memory reclaim, page demotion, or NUMA placement can be the difference between stable tail latency and service-level objective violations. At the same time, operators want policies that are specialized to their own services, hardware, and workload phases. This has made kernel programmability attractive: eBPF programs can encode service-specific logic and run inside privileged kernel paths after verifier checks.

This trend is moving beyond packet filters and tracing. sched\_ext allows scheduling behavior to be defined by BPF programs. PageFlex-style mechanisms show that paging decisions can also be delegated or customized through BPF-assisted policy paths. These mechanisms let operators specialize scheduling and memory policy for networked services, but they also introduce a new failure mode: independently safe policies may create harmful resource effects when composed.

Consider a latency-sensitive key-value service under memory pressure. A BPF scheduler boosts the service when queue delay increases. A BPF paging policy, using stale phase information, demotes pages from the same service. Each policy is reasonable alone: the scheduler reduces queueing, and the paging policy reduces memory footprint. Together, however, they can form a feedback loop. The scheduler gives the service more execution opportunities; the service touches pages that the paging policy demoted; refaults increase; queue delay and tail latency rise; the scheduler boosts again; the paging policy continues demoting. The result is a verifier-accepted but harmful composition.

The key observation is that verifier acceptance is not resource-effect safety. The verifier can reject unsafe memory accesses or unbounded control flow, but it does not prove that a policy creates only a bounded amount of queueing, demotion, refault, or cross-subsystem debt for a service scope. Nor does it specify how a system should recover when a policy remains memory-safe but becomes effect-unsafe.

We propose \emph{resource-effect safety} as a missing safety dimension for programmable kernel policy. Resource-effect safety asks three questions. First, which resource effects may a policy create? Second, how much resource debt may each effect create within a service scope? Third, when a budget is violated, can the system recover by revoking only the harmful effect rather than disabling all BPF functionality?

\sys answers these questions using three mechanisms. First, \emph{effect tokens} authorize a BPF policy to create specific effects, such as \texttt{boost\_task} or \texttt{demote\_page}, within a scope and budget. Second, \emph{per-scope resource ledgers} charge scheduler and paging effects to shared service scopes, connecting cgroup-level scheduling events and memcg-level memory events. Third, \emph{bounded degradation} throttles or revokes the smallest harmful effect that explains a violation.

This paper targets NSDI because the motivating problem is not merely kernel extensibility. It is the stability of networked services whose resource-management policies are increasingly programmable. The current artifact combines controlled QEMU validation, QEMU memcached matrices, and native non-QEMU VM memcached P5/P6 bars. A submission-ready evaluation still needs a broader bare-metal or CXL/PMEM campaign across latency-sensitive services, multi-tenant pressure, and service-level recovery.

\paragraph{Contributions.}
This paper makes four contributions:
\begin{itemize}[leftmargin=*]
    \item We identify resource-effect conflicts among verifier-accepted BPF policies as a safety problem for programmable resource management in networked services.
    \item We define effect tokens and per-scope ledgers for accounting scheduling and paging effects across shared service scopes.
    \item We design effect-boundary enforcement and bounded degradation, allowing the system to revoke harmful effects while preserving unrelated safe policy functionality.
    \item We implement \sys on Linux using sched\_ext and a conservative PageFlex-style paging decision hook, and validate it with QEMU scheduler-paging conflicts plus a memcached service matrix.
\end{itemize}

\section{Background and Motivation}

\subsection{BPF safety today}

The eBPF verifier is a gatekeeper for kernel-resident BPF programs. It reasons about program control flow, bounded execution, register and stack state, pointer types, helper arguments, and memory accesses. This prevents many classes of unsafe programs from entering the kernel. However, verifier acceptance does not imply that the policy encoded by a program is benign under every runtime condition. A scheduler policy can legally reorder tasks in a way that increases starvation risk. A paging policy can legally request demotion decisions that later increase refaults. These are policy-effect failures rather than verifier failures.

\subsection{Programmable scheduling}

sched\_ext exposes a BPF-defined scheduler interface. A BPF scheduler can decide how tasks are enqueued, dispatched, and selected for CPUs. This makes it possible to specialize scheduling policy for service-level objectives, but it also allows policies to create resource effects such as boosting, pinning, dispatch failures, and queue-delay shifts.

\subsection{Programmable paging}

PageFlex-style systems show that paging policy decisions can be delegated while exposing kernel memory state at low overhead. This creates a complementary source of resource effects: demotion, reclaim hints, region classification, refault amplification, and fault-latency changes.

\subsection{The missing cross-subsystem layer}

The scheduler and memory manager are coupled through workload execution. A scheduling decision changes which pages are touched. A paging decision changes how quickly tasks can make progress. If both sides become programmable, local safety is insufficient. We need a way to account for cross-subsystem effects at the service scope where the application experiences them.

\section{Problem Statement}

\subsection{Threat model}

We assume BPF programs pass the Linux verifier. We do not address verifier soundness bugs, JIT bugs, arbitrary kernel memory corruption, or hardware side channels. We focus on policies that are locally legal but may be malicious, misconfigured, stale, or generated by independent controllers with conflicting objectives. Policies may belong to different teams or tenants.

\subsection{Resource-effect conflict}

A resource-effect conflict occurs when two or more legal policy effects jointly violate a runtime resource invariant. In this paper we focus on scheduler-paging conflicts:

\begin{quote}
A scheduler effect increases execution opportunities for a service scope while a paging effect increases refault or fault-latency debt for the same scope, causing queue-delay and tail-latency amplification.
\end{quote}

\subsection{Design goals}

\begin{enumerate}[leftmargin=*]
    \item \textbf{Post-verifier enforcement.} Do not replace the verifier.
    \item \textbf{Effect-level attribution.} Attribute debt to effects, not just programs.
    \item \textbf{Cross-subsystem scopes.} Account scheduler and paging effects in the same service scope.
    \item \textbf{Low overhead.} Enforce at effect boundaries, not every instruction.
    \item \textbf{Bounded degradation.} Prefer throttling or revoking harmful effects over killing the whole policy.
\end{enumerate}

\section{Design}

\subsection{Overview}

\sys consists of a user-space contract manager and a kernel resource-effect layer. The contract manager parses manifests, installs effect tokens, resolves scopes, and exports audit events. The kernel layer stores token tables, updates per-scope ledgers, checks effect gates, and triggers degradation.

\begin{table}[t]
\centering
\caption{Examples of resource effects.}
\begin{tabular}{lll}
\toprule
Subsystem & Effect & Resource debt \\
\midrule
sched\_ext & boost\_task & queue delay, boost rate \\
sched\_ext & dispatch\_task & dispatch failures \\
sched\_ext & pin\_cpu & CPU imbalance, starvation \\
Paging & demote\_page & refaults, fault latency \\
Paging & reclaim\_hint & reclaim pressure \\
Paging & classify\_region & stale phase decisions \\
\bottomrule
\end{tabular}
\end{table}

\subsection{Effect tokens}

An effect token authorizes a resource effect within a scope and budget:
\[
Token = \langle subsystem, effect, scope, budget, fallback \rangle.
\]

For example, a paging policy may receive a token authorizing
\texttt{demote\_page} for memcg A with a refault budget and a
\texttt{revoke\_demote} fallback.

This differs from helper permissions. A helper allowlist asks whether a program may call an API. An effect token asks whether a policy may create a bounded resource effect for a given service scope.

\subsection{Per-scope resource ledgers}

A ledger aggregates resource debt for a scope:

\begin{lstlisting}[language=C,basicstyle=\ttfamily\small]
struct ledger {
  scope_id scope;
  u64 sched_queue_delay_us;
  u64 sched_dispatch_failures;
  u64 sched_boost_events;
  u64 pages_demoted;
  u64 refault_events;
  u64 major_fault_events;
  u64 fault_latency_us;
  u32 degrade_state[EFFECT_MAX];
};
\end{lstlisting}

The scope can be a cgroup, memory cgroup, CPU set, NUMA node, or memory region. The paper focuses on service scopes that connect cgroup and memcg identities.

\subsection{Effect-boundary enforcement}

\sys checks only when a BPF policy attempts to create a resource effect. For sched\_ext, boundaries include select\_cpu, enqueue, dispatch, boost, and CPU pinning. For paging, boundaries include demotion decisions, reclaim hints, and region classification.

\begin{lstlisting}[language=C,basicstyle=\ttfamily\footnotesize]
int gate(prog, effect, scope, cost) {
  token = lookup_token(prog, effect, scope);
  if (!token) return DENY;

  ledger = lookup_ledger(scope);
  if (would_exceed(ledger, token, cost)) {
    trigger_degrade(prog, effect, scope, ledger);
    return DEGRADED;
  }

  charge(ledger, effect, cost);
  return ALLOW;
}
\end{lstlisting}

\subsection{Cross-subsystem rule}

The minimal rule used in the prototype is:

\begin{quote}
If a service scope has high refault ratio, high queue delay, and high demotion rate in the same epoch window, revoke the demote\_page effect for that scope while preserving scheduler dispatch.
\end{quote}

This rule is intentionally simple. The paper's claim is not that this is the only possible invariant, but that a small set of resource-effect invariants can prevent damaging feedback loops.

\subsection{Bounded degradation}

\begin{table}[t]
\centering
\caption{Bounded degradation levels.}
\begin{tabular}{lll}
\toprule
Level & Action & Example \\
\midrule
0 & normal execution & allow boost and demotion \\
1 & throttle effect & reduce demotion rate \\
2 & revoke effect & revoke demote\_page \\
3 & subsystem fallback & default kernel reclaim \\
4 & disable policy & disable sched\_ext policy \\
\bottomrule
\end{tabular}
\end{table}

The key design choice is minimality. If refault debt is caused by page demotion, \sys revokes demotion before disabling the scheduler. This preserves useful BPF functionality and reduces disruption.

\section{Implementation}

We implement \sys in Linux using three components: a contract manager, a sched\_ext front-end, and a paging front-end.

The artifact pins Linux 6.12.30 and keeps the kernel changes as a patch series
under \texttt{kernel/patches/}. The current series contains 2,294 patch lines
across the core token and ledger layer, sched\_ext gate, conservative MM
decision hook, cross-subsystem rule, and kselftest wiring. The repository also
contains QEMU rootfs builders, run scripts, BPF policy stubs, Rust user-space
tools, workload drivers, and experiment parsers. Generated Linux trees,
initramfs images, raw logs, and figures are kept out of version control.

\subsection{Contract manager}

The user-space manager parses YAML manifests, validates effect names and fallback availability, resolves scopes, and installs token maps. It also streams violation events for evaluation.

\subsection{sched\_ext integration}

The current sched\_ext front-end gates the dispatch/boost path used by the
prototype scheduler. It records queue delay, boost rate, and dispatch activity
using per-scope counters. The broader design also targets select\_cpu, enqueue,
and CPU-pinning boundaries, but those are not yet all implemented in the
artifact.

\subsection{Paging integration}

The paging front-end follows a conservative decision-only design. BPF programs receive summarized read-only state and return keep, demote, reclaim\_hint, or no\_op. The kernel validates page state and effect budgets before execution.

\subsection{Overhead controls}

\sys uses per-CPU counters, epoch aggregation, static token lookup caching, and optional sampling for high-rate paging events. These choices are necessary to keep enforcement overhead low.

\section{Prototype Evaluation}

This section reports the evidence currently produced by the artifact. All
experiments below run in QEMU on the pinned Linux 6.12.30 kernel. We use QEMU to
validate correctness and reproducibility of the kernel paths; the absolute
latency and throughput numbers are not production performance results.

\subsection{Workloads}

The main matrix uses a controlled phase-changing service that touches memory in
phases and emits per-phase latency samples. This workload is intentionally small
so that every baseline can run in one QEMU boot and produce comparable raw
serial evidence. The artifact also includes a real memcached path:
\texttt{make memcached-experiments} boots QEMU, starts two memcached instances
in the guest, and runs an ASCII protocol load client across the same G1--G9
baseline set. This matrix uses controlled kernel effect injection for the
scheduler-paging conflict; it is a real-service artifact validation, not a
production performance experiment.

\subsection{Baselines}

\begin{table}[t]
\centering
\small
\caption{Evaluation baselines.}
\begin{tabular}{ll}
\toprule
Group & Configuration \\
\midrule
G1 & default Linux scheduler and paging \\
G2 & sched\_ext policy only \\
G3 & BPF paging policy only \\
G4 & sched\_ext + BPF paging without \sys \\
G5 & cgroup/memcg quota-style controls \\
G6 & static checker only \\
G7 & per-subsystem ledger only \\
G8 & kill-whole-policy fallback \\
G9 & full \sys \\
\bottomrule
\end{tabular}
\end{table}

\subsection{Metrics}

We measure P50/P99/P999 latency, throughput, queue delay, boost events, pages demoted per epoch, refault ratio, major fault rate, fault latency, fallback activation latency, recovery time, and steady-state overhead.

\subsection{Controlled QEMU matrix}

% Generated by experiments/analysis/generate_paper_tables.py.
% Source: experiments/results/processed/matrix_summary.csv
% Source-SHA256: 2c01bcbea1d19950817281b3ce4aa0fa25938e565394820dbfa524fee1924db4

\begin{table}[t]
\centering
\scriptsize
\caption{Controlled QEMU matrix summary. Latencies are guest synthetic-service samples.}
\begin{tabular}{lrrrr}
\toprule
Group & P99 us & Queue delay us & Demoted & Demote state \\
\midrule
G1 & 95243 & 0 & 0 & 0 \\
G2 & 95088 & 0 & 0 & 0 \\
G3 & 93203 & 0 & 8 & 0 \\
G4 & 93092 & 0 & 0 & 0 \\
G5 & 94300 & 0 & 0 & 0 \\
G6 & 95674 & 0 & 0 & 0 \\
G7 & 96894 & 36000 & 8 & 0 \\
G8 & 92413 & 34000 & 8 & 0 \\
G9 & 93680 & 31000 & 8 & 2 \\
\bottomrule
\end{tabular}
\end{table}


The matrix produces raw logs under \texttt{experiments/results/raw/}, processed
CSV tables under \texttt{experiments/results/processed/}, and five SVG figures:
feedback timeline, tail latency, recovery, ablation, and overhead. This
controlled matrix is primarily an artifact plumbing and baseline-validation
table; the latest naturally induced conflict and recovery evidence is reported
separately below. The recovery table still records the \sys path reaching
demote degrade state 2, which corresponds to revoking demote\_page, while
preserving scheduler dispatch.

\subsection{Real-service companion matrix}

% Generated by experiments/analysis/generate_paper_tables.py.
% Source: experiments/results/processed/memcached_matrix_summary.csv
% Source-SHA256: 200df92b6a0943f1e813aff9d03832aeea1f431d4941388445a5451269298837

\begin{table}[t]
\centering
\scriptsize
\caption{QEMU memcached companion matrix.}
\begin{tabular}{lrrrrr}
\toprule
Group & P99 us & Tenant-B P99 & Queue us & Demoted & Demote state \\
\midrule
G1 & 5119 & 5965 & 0 & 0 & 0 \\
G2 & 3996 & 3873 & 0 & 0 & 0 \\
G3 & 4174 & 4314 & 0 & 8 & 0 \\
G4 & 4011 & 2030 & 0 & 0 & 0 \\
G5 & 696 & 1183 & 0 & 0 & 0 \\
G6 & 977 & 982 & 0 & 0 & 0 \\
G7 & 4115 & 4041 & 106000 & 8 & 2 \\
G8 & 18414 & 4303 & 89000 & 8 & 0 \\
G9 & 5881 & 2630 & 155000 & 8 & 2 \\
\bottomrule
\end{tabular}
\end{table}


The memcached matrix proves that the artifact can boot a guest, start real
networked services, drive load, and collect latency samples across every
baseline group. The raw log is
the May 22 memcached-matrix QEMU log, and the processed table is
\texttt{memcached\_matrix\_summary.csv}. This companion matrix is not the final
conflict table; it establishes that the real-service QEMU harness can run every
baseline group and produce auditable latency samples.

% Generated by experiments/analysis/generate_paper_tables.py.
% Source: experiments/results/processed/memcached_natural_bars.csv
% Source-SHA256: 824233f2c89620dc452d8de5f340a6406a166b4ec7548f534a5e281f6541abb9

\begin{table}[t]
\centering
\scriptsize
\caption{QEMU memcached natural conflict and recovery bars.}
\begin{tabular}{lrrrrr}
\toprule
Group & P99 us & Tenant-B P99 & Queue us & Refaults & Demote state \\
\midrule
G1 & 3472 & 16467 & 0 & 0 & 0 \\
G2 & 3385 & 3765 & 0 & 0 & 0 \\
G4 & 402778 & 101720 & 3852000 & 28344 & 2 \\
G9 & 3312 & 3469 & 863000 & 28328 & 2 \\
\bottomrule
\end{tabular}
\end{table}

\IfFileExists{generated/native_memcached_bars_table.tex}{%
% Generated by experiments/analysis/generate_paper_tables.py.
% Source: experiments/results/processed/native_memcached_bars.csv
% Source-SHA256: ff9d32659695b61dc29d5d76cb6c098da90beb81c10efe6ff16afbcd45f35e12

\begin{table}[t]
\centering
\scriptsize
\caption{Native non-QEMU memcached conflict and recovery bars.}
\begin{tabular}{lrrrrr}
\toprule
Group & P99 us & Tenant-B P99 & Queue us & Refaults & Demote state \\
\midrule
G1 & 8202 & 378 & 0 & 0 & 0 \\
G2 & 8354 & 137 & 0 & 0 & 0 \\
G4 & 167208 & 75 & 225184000 & 186268 & 2 \\
G9 & 18455 & 63 & 433953000 & 251475 & 2 \\
\bottomrule
\end{tabular}
\end{table}

Native non-QEMU memcached conflict and recovery bars are generated from the
native final-evidence CSV and supersede the QEMU proxy for P5/P6 performance
claims.
}{}

The natural-bars memcached run removes the controlled counter injection used by
the companion matrix and instead drives real memory pressure in the QEMU guest.
G4 naturally raises demotion, refault, and queue-delay counters, while G9
reaches demote degrade state 2 and records a lower recovery-window P99. The
native table above provides the current non-QEMU VM P5/P6 validation; the paper
keeps QEMU and native VM evidence separate from bare-metal production claims.

% Generated by experiments/analysis/generate_paper_tables.py.
% Source: experiments/results/processed/no_violation_overhead.csv
% Source-SHA256: 1fa241066f7199b293fee7f34c1a9460156f0146c238cbc09cdc76f5e6aa1526

\begin{table}[t]
\centering
\scriptsize
\caption{QEMU no-violation overhead probe.}
\begin{tabular}{lrrr}
\toprule
Phase & P99 overhead \% & Throughput overhead \% & CPU overhead \% \\
\midrule
NV1 & -90.44 & -130.22 & -2.93 \\
\bottomrule
\end{tabular}
\end{table}


The no-violation probe checks the current QEMU artifact's steady-state overhead
when ContractBPF gates are installed but the cross-subsystem violation is not
triggered. Production overhead claims still require non-QEMU runs.

\subsection{Current findings}

\begin{enumerate}[leftmargin=*]
    \item The kernel patch series builds, QEMU boots, sched\_ext loads and unloads, contractd starts, and kselftests exercise token, ledger, degrade, MM hook, and cross-subsystem paths.
    \item The natural QEMU memory-pressure scenario records the feedback-loop counters needed for the thesis: queue delay, demotion, and refault or major-fault activity without \texttt{cross\_scenario} counter injection.
    \item The full \sys path detects the cross-subsystem condition and revokes only demote\_page while keeping scheduler dispatch active.
    \item A real memcached service can run in QEMU across the full G1--G9 matrix, and the native non-QEMU VM run records same-load P5/P6 conflict and recovery bars.
\end{enumerate}

\section{Discussion}

\subsection{Why not cgroups?}

cgroups bound coarse resource usage but do not attribute debt to specific BPF effects. A cgroup CPU quota can throttle service execution, but it cannot revoke only a demote\_page effect while preserving safe scheduler dispatch.

\subsection{Why not static checking?}

Static checks can reject obvious manifest conflicts, but the motivating failure depends on runtime phase, memory pressure, refaults, and queue delay. These are dynamic effects.

\subsection{Why NSDI rather than a pure OS venue?}

The mechanisms are kernel mechanisms, but the failure target is networked services. Programmable resource policies are becoming part of cloud service management, so a strong NSDI submission must emphasize multi-tenant service stability, tail latency, workload phases, and operational usefulness. The current artifact provides kernel, QEMU, and native VM evidence for that argument, but it still needs a broader bare-metal real-service evaluation.

\subsection{Limitations}

\sys currently targets sched\_ext and paging. It does not handle all BPF program types, does not replace the verifier, and does not prove global optimality of resource policies. Threshold selection and scope mapping require careful configuration.

The current evaluation is also limited. QEMU validates kernel control flow,
reproducibility, and failure recovery, but it is not evidence of production
performance. The native P5/P6 bars run on a VMware VM with an explicit
ContractBPF lower-tier node model, so they provide non-QEMU final-evidence
coverage for these gates without claiming bare-metal CXL/PMEM behavior. The
paging path is a conservative kernel-validated prototype, not the final
PageFlex-style BPF loading path. The paper must not claim submission-ready
production performance until dedicated bare-metal multi-tenant experiments are
added.

\section{Related Work}

\paragraph{eBPF verifier.}
The verifier provides local execution safety for BPF programs \cite{linuxVerifier}. \sys is complementary: it handles runtime resource effects after verifier acceptance.

\paragraph{sched\_ext.}
sched\_ext enables BPF-defined scheduling policies \cite{linuxSchedExt}. \sys adds bounded effect accounting for scheduling decisions.

\paragraph{BPF-assisted paging.}
PageFlex demonstrates flexible paging policy delegation using eBPF-style mechanisms \cite{pageflex2025}. \sys studies what happens when such paging policies interact with programmable scheduling.

\paragraph{Static eBPF policy enforcement.}
KRAKENGUARD enforces policy-driven constraints on BPF bytecode at load time and detects cross-program interference \cite{krakenguard2026}. \sys instead targets runtime state-dependent resource-effect debt.

\paragraph{Safe kernel extensions.}
Rex addresses language and verifier usability for safe kernel extensions \cite{rex2025}. \sys addresses resource-effect conflicts even when code is safe.

\section{Conclusion}

BPF programmability is moving into core kernel policy paths that affect networked services. Local verifier acceptance is necessary but not sufficient when independently loaded policies create interacting resource effects. \sys introduces effect tokens, per-scope ledgers, and bounded degradation to account for and recover from scheduler-paging resource-effect conflicts. The central lesson is that programmable kernel policy needs runtime effect safety, not just program safety.

\bibliographystyle{plain}
\bibliography{references}

\end{document}
